\chapter{Conclusión}

El desarrollo de SmartGate permitió cumplir con el objetivo principal del proyecto: implementar un sistema de control de acceso mediante reconocimiento facial, capaz de realizar el procesamiento y la toma de decisiones de forma local sobre una Raspberry Pi. La integración de los modelos de detección facial, antispoofing y reconocimiento, junto con la electrónica de control del portón, permitió obtener un prototipo funcional que reúne las características definidas para el producto mínimo viable (MVP).

Uno de los principales resultados fue demostrar la viabilidad de ejecutar múltiples modelos de machine learning sobre una plataforma embebida de recursos limitados, sin recurrir a aceleradores especializados ni depender de servicios externos para realizar las inferencias. La arquitectura modular y el procesamiento basado en eventos permitieron coordinar las distintas etapas, administrar los recursos disponibles y mantener un funcionamiento ordenado frente a cambios de estado del portón o errores durante la ejecución.

Además del cumplimiento del MVP, se logró incorporar gran parte de las funcionalidades complementarias previstas. La integración con Supabase, la interfaz web de administración, la persistencia local y los mecanismos de sincronización permitieron ampliar considerablemente las capacidades del sistema. Estas implementaciones aportaron autonomía, trazabilidad y facilidad de administración, manteniendo la posibilidad de operar ante interrupciones temporales de conectividad.

Las pruebas realizadas con distintas personas permitieron comprobar el funcionamiento conjunto de los módulos y obtener resultados favorables para esta etapa de desarrollo. Sin embargo, también se identificaron limitaciones relacionadas principalmente con las condiciones de captura de imágenes y la capacidad de generalización de los modelos utilizados. El entrenamiento específico con un dataset propio no pudo completarse debido al tiempo disponible y a las dificultades encontradas durante la adquisición de muestras, aunque se desarrollaron herramientas que permitirán retomar esta tarea posteriormente.

Desde el punto de vista de la planificación, las principales etapas y entregas pudieron mantenerse dentro de los plazos establecidos, adaptando las prioridades frente a las dificultades de disponibilidad de hardware y a la complejidad de determinadas funcionalidades. Esta experiencia evidenció la importancia de establecer objetivos alcanzables, desarrollar de forma incremental y validar progresivamente cada módulo antes de avanzar hacia una integración completa.

En definitiva, SmartGate constituye una implementación funcional que demuestra la posibilidad de combinar visión artificial, sistemas embebidos, diseño electrónico y servicios de administración remota en una solución de control de acceso. Si bien todavía requiere una validación más extensa y mejoras para su utilización en un entorno real de manera permanente, el desarrollo alcanzado establece una base técnica sólida y ampliable, sobre la cual es posible continuar optimizando el reconocimiento, la seguridad y las prestaciones generales del sistema.
